% Report Module 6 – Industrial Espionage in Cyberspace
% CS 475: Introduction to Computer Security
% Mária Nyolcas
% March 2026

\documentclass[11pt,a4paper]{article}

\usepackage[margin=2.5cm]{geometry}
\usepackage{setspace}
\usepackage{booktabs}
\usepackage{tabularx}
\usepackage{array}
\usepackage{parskip}
\usepackage{titlesec}
\usepackage{hyperref}
\usepackage{enumitem}
\usepackage[protrusion=true,expansion=false]{microtype}
\usepackage{fancyhdr}
\usepackage{graphicx}
\usepackage{xcolor}

\setstretch{1.15}

\hypersetup{
    colorlinks=true,
    linkcolor=black,
    urlcolor=blue,
    citecolor=black,
}

\titleformat{\section}{\large\bfseries}{}{0em}{}[\titlerule]
\titleformat{\subsection}{\normalsize\bfseries}{}{0em}{}

\pagestyle{fancy}
\fancyhf{}
\lhead{\small CS 475 – Mária Nyolcas}
\rhead{\small Week 6: Industrial Espionage in Cyberspace}
\cfoot{\thepage}
\renewcommand{\headrulewidth}{0.4pt}

\begin{document}

\begin{center}
  {\Large\bfseries Report Module 6}\\[4pt]
  {\large Industrial Espionage in Cyberspace}\\[6pt]
  {\normalsize CS 475: Introduction to Computer Security}\\[2pt]
  {\normalsize Mária Nyolcas \quad|\quad Prof: Asen Nikolov Grozdanov}\\[2pt]
  {\normalsize March 2026}
\end{center}

\vspace{4pt}
\noindent\rule{\linewidth}{0.6pt}
\vspace{2pt}

\noindent\textbf{Abstract.}
This report documents a structured OSINT (Open-Source Intelligence) reconnaissance
investigation of a fictional target organisation, InnovateTech Solutions, alongside an
analysis of the five-phase AeroTech GmbH espionage lifecycle case study.  The investigation
demonstrates how publicly available information can be assembled into an actionable threat
intelligence profile without touching the target's systems.  Defensive countermeasures are
proposed for each phase, the discussion questions from Section~6 of the module brief are
answered with reference to the lab evidence, and reflections connect the findings to the
broader course arc.


\section{Introduction}

Industrial espionage --- the covert acquisition of confidential business or technical
information for competitive advantage --- has shifted almost entirely into cyberspace.
What once required a bribed insider or a stolen briefcase can now be accomplished remotely,
at scale, with a level of plausible deniability that physical theft never allowed.  ENISA's
annual threat landscape consistently ranks cyber-enabled economic espionage among the
top threats facing European organisations, and the FBI has described it as ``the greatest
long-term threat to the nation's economy'' \cite{fbi2023}.

For a computer science student, this topic is not abstract.  The systems we design and the
code we write will one day protect --- or inadvertently expose --- trade secrets, R\&D data,
and competitive intelligence.  Understanding how adversaries operate, and how little
information they need to build a dangerous profile, is the first step toward building
defences that actually work.

This report covers: (i) the methodology and findings of a passive OSINT investigation of
the fictional target InnovateTech Solutions; (ii) analysis of the AeroTech GmbH composite
case study across all five attack phases; (iii) defensive recommendations grounded in NIST,
CIS Controls, and ISO/IEC 27001; and (iv) reflections on how this module connects to
earlier weeks and the road ahead.


\section{Methodology}

The OSINT investigation followed the six-step structure prescribed in the lab exercise.
All findings relate to the fictional company InnovateTech Solutions
(\texttt{innovatech-solutions.io}, advanced materials / aerospace R\&D, Munich).
No real organisation was investigated.  The investigation was entirely passive: no active
scanning, no network contact with any system.

\subsection{Tools and Sources Used}

\begin{center}
\small
\begin{tabular}{lll}
\toprule
Step & Tool / Source & Purpose \\
\midrule
Passive web recon  & Search engine (advanced operators), WHOIS & Domain, registrar, admin contacts \\
Legacy content     & Wayback Machine                           & Historical pages, old tech references \\
Code repositories  & GitHub search                             & Accidentally committed secrets \\
Patent intelligence & EPO patent database                      & R\&D direction, researcher names \\
Employee profiling & LinkedIn, ResearchGate                    & Roles, skills, project names \\
Technical footprint & DNS (\texttt{dig}), crt.sh               & Subdomains, exposed services \\
\bottomrule
\end{tabular}
\end{center}

\vspace{4pt}
\noindent
All tool interactions were additionally encoded in \texttt{osint\_recon.py}, a structured
Python script that stores findings as typed dataclasses, making the evidence reproducible
and queryable --- following the same design philosophy as previous weeks' lab scripts.


\section{Findings}

\subsection{Passive Web Reconnaissance}

Five high-value findings emerged from web sources alone.  The most significant was a
\texttt{config.yml} file pushed to a public GitHub repository (\texttt{innovatech-dev})
by a junior engineer in November 2022.  It contained the internal subnet
\texttt{10.50.0.0/16} and the staging hostname \texttt{sim-lab01.internal} --- information
that would allow an adversary to target internal reconnaissance without any active
scanning of the external perimeter.

A 2021 Wayback Machine snapshot of the careers page disclosed the VPN product and version
(\textit{Cisco AnyConnect 4.9}), which is affected by CVE-2021-1247.  Job postings named
specific ANSYS and GitLab CE versions, directly mapping to CVE lookup targets.

\subsection{Employee and Personnel Profiling}

Four persons of interest were identified.  Dr.~K.~Brenner (Head of Materials Research)
listed the internal project name \textit{ComposiX-II} in their LinkedIn experience section.
T.~Okafor (Senior IT Administrator) had contributions to \texttt{innovatech-dev/infra-scripts}
with a commit referencing a Vault token path.  Dr.~A.~Reyes is a co-inventor on three
recent patents, with a conference biography naming the active research lab location.

The junior engineer M.~Hofer, whose email is visible in the git log of the leaking commit,
represents the most accessible initial access target: already demonstrably careless with
internal data and likely to respond to a convincing spear-phishing pretext.

\subsection{Technical Footprinting}

DNS enumeration revealed that the self-hosted GitLab instance
(\texttt{gitlab.innovatech-solutions.io}) is directly internet-exposed with no IP
restriction.  The SPF record uses \texttt{\textasciitilde{}all} (softfail) rather than
\texttt{-all}, meaning spoofed sender addresses from the \texttt{innovatech-solutions.io}
domain will reach most mail servers without rejection.  No DMARC record was found.
Certificate transparency logs (crt.sh) disclosed four subdomains including
\texttt{sim-lab01} and \texttt{dev-api}, both of which should not be visible in public
certificate logs.

\subsection{Synthesised Attack Surface Summary}

\small
\begin{tabularx}{\linewidth}{l >{\raggedright\arraybackslash}X >{\raggedright\arraybackslash}X}
\toprule
Finding & Adversarial value & Remediation \\
\midrule
Leaked subnet + hostname (GitHub) & Internal recon without network contact & Secrets scanning pre-commit hook (gitleaks) \\
No DMARC / SPF softfail           & Spoofed sender phishing at scale      & Add DMARC p=reject; harden SPF to -all \\
GitLab internet-exposed           & Unauthenticated code access           & Move behind VPN / IP allowlist \\
VPN version disclosed (Wayback)   & CVE-targeted exploit                  & Remove version strings from public content \\
Researcher named on patents       & Targeted spear-phishing pretext       & Omit project names from public bios \\
\bottomrule
\end{tabularx}
\normalsize


\section{Analysis}

\subsection{AeroTech GmbH Case Study}

The five-phase lifecycle maps cleanly onto the Lockheed Martin Cyber Kill Chain.
Phase~1 (Reconnaissance) and Phase~2 (Initial Access) follow the same pattern as
InnovateTech: passive OSINT enabling a highly credible spear-phish.  The critical
observation is that \textbf{Phase~2 is the single best disruption point} (Discussion
Question~1).  Once the PDF exploit executes, the attacker has a foothold and every
subsequent phase compounds the advantage.  Blocking delivery --- through email sandboxing,
mandatory patching of document viewers, and phishing simulation training --- collapses the
chain before it starts.  Hardening Phase~3 or~4 is valuable but reactive: the adversary
is already inside.

\textbf{Least privilege and Phase~3} (Discussion Question~2): the lateral movement
succeeded because shared service accounts had broad access across network segments.
If the procurement workstation had been separated from the R\&D file server by a
firewall requiring explicit authorisation, the attacker would have stalled at the first
segment boundary.  The blast radius of the initial compromise would have been
limited to a single low-value host.  This directly echoes the Bell-LaPadula ``no read up''
principle from Week~2: mandatory, non-discretionary boundaries contain damage even when
perimeter defences fail.

\textbf{Human behaviour} (Discussion Question~3): no phase was purely technical.
Phase~1 required a junior engineer to commit secrets; Phase~2 required a researcher to
open an attachment; Phase~3 depended on administrators sharing accounts for convenience;
Phase~5 relied on a single analyst reviewing historical logs during a routine audit.
Technology failed at every step because it was misconfigured or absent --- the root cause
in each case was a human decision, not an undefeatable technical capability.

\subsection{DLP Policy Requirements}

A DLP policy addressing the risks identified would need three layers.
\textit{At the endpoint}: block bulk copy of files tagged as R\&D or Confidential to
removable media or personal cloud storage.
\textit{At the network boundary}: alert on outbound transfers exceeding a volume threshold
to cloud storage domains not on an approved list; inspect TLS for anomalous certificate
chains.
\textit{At the repository}: pre-receive hooks rejecting commits containing patterns matching
IP address ranges, hostnames, or secret tokens (gitleaks or truffleHog).

\subsection{Highest Impact-to-Cost Defensive Control}

Across all findings, the single control with the highest impact-to-cost ratio is
\textbf{secrets scanning in the CI/CD pipeline}.  A pre-commit hook running gitleaks
costs near zero to deploy, requires no ongoing licensing, and would have prevented the
most actionable finding in this investigation --- the leaked subnet and hostname --- before
it ever reached a public repository.  Every other finding derived some portion of its
value from that one commit.


\section*{Conclusion}

This module demonstrated that the most dangerous phase of a cyber espionage campaign
requires no hacking at all: passive OSINT alone produced five high-value findings against
a fictional target, including internal network topology, software CVE targets, and direct
candidate identities for spear-phishing.  The AeroTech case study showed that human
decisions --- a careless commit, a clicked attachment, a shared account --- created every
exploitable gap, and that the highest-leverage defensive investment is therefore in
people and process, not exclusively in technology.

Looking ahead, the coming modules on operating systems security and network defences will
address how to harden the environment that an attacker enters \textit{after} initial
access.  Industrial espionage makes clear why those controls matter: an attacker who spends
four weeks moving laterally undetected has effectively unlimited time to find and extract
the crown jewels.  The goal of layered defence is to ensure that no single failure ---
human or technical --- is sufficient to reach them.


\section*{References}

\begin{enumerate}[label={[\arabic*]}, leftmargin=*, nosep]
  \item MITRE Corporation. (2024). \textit{ATT\&CK Enterprise Matrix v15}. \url{https://attack.mitre.org}
  \item ENISA. (2024). \textit{ENISA Threat Landscape 2024}. \url{https://www.enisa.europa.eu/publications/enisa-threat-landscape}
  \item National Institute of Standards and Technology. (2018). \textit{Framework for Improving Critical Infrastructure Cybersecurity v1.1}. \url{https://www.nist.gov/cyberframework}
  \item Lockheed Martin. (2022). \textit{Cyber Kill Chain}. \url{https://www.lockheedmartin.com/en-us/capabilities/cyber/cyber-kill-chain.html}
  \item Center for Internet Security. (2024). \textit{CIS Controls v8}. \url{https://www.cisecurity.org/controls}
  \item Council of Europe. (2001). \textit{Budapest Convention on Cybercrime}. \url{https://www.coe.int/en/web/cybercrime/the-budapest-convention}
  \item European Parliament. (2016). \textit{Directive 2016/943 on the protection of undisclosed know-how and business information (trade secrets)}. \url{https://eur-lex.europa.eu}
  \item IP Commission. (2017). \textit{The Theft of American Intellectual Property: Reassessments of the Challenge and United States Policy}. \url{https://www.ipcommission.org}
  \item\label{fbi2023} Federal Bureau of Investigation. (2023). \textit{Economic Espionage}. \url{https://www.fbi.gov/investigate/counterintelligence/economic-espionage}
  \item SANS Institute. (2024). \textit{Reading Room: Insider Threat and Data Exfiltration}. \url{https://www.sans.org/reading-room/}
\end{enumerate}

\end{document}
